\section{Discussion on metaprogramming}
\label{discussion}

\subsection{Macro expansion}
\label{macroexpand-mechanism}
The \(\ildsf{macroexpand}\) special form allows metaprogramming by treating a certain function
as a \emph{macro}. A regular function evaluation \(\ildlist{f\, a_1\, a_2\, \dots\, a_n}\) evaluates f and all
arguments and then passes the evaluated arguments to the evaluated f. In contrast,
\(\ildlist{\ildsf{macroexpand}\, f\, a_1\, a_2\, \dots\, a_n}\) evaluates just \(f\) and then passes the
\textbf{unevaluated} arguments to it. The result is then in turn evaluated. This allows \(f\)
to treat the program passed to it as data and to transform it arbitrarily before it
gets evaluated. This is similar to unhygienic macro systems like the one in LISP.
As in traditional LISP macro systems \cite{kohlbecker1986syntactic}, ILD macros
expand from the outside in: the outermost macro call is expanded first, and its
expansion may itself contain further macro calls.

Unlike LISP, ILD denotes macro expansion at the callsite rather than
differentiating between \emph{functions} and \emph{macros}. This greatly simplifies the
semantics and implementation.

Another difference from LISP macros is that ILD does not have a separate macro
expansion phase: instead, macros are evaluated as encountered. We will call
this \emph{runtime semantics of macro expansion}. This lack of separate macro expansion
phase is similar to the way metaprogramming works in FExpr-based languages\cite{fexpr-shutt},
and has advantages (for example, generated code can depend on runtime input).
Nevertheless, performance can be greatly improved by evaluating macros ahead of time.
We argue that placing the boundary between code evaluated ahead of time and code
evaluated at runtime should not be at the distinction between ``macro expansion'' and
``code execution'', but rather between ``code dependant on data known ahead of time''
(which includes most, but not all, macros) and ``code dependant on runtime data''.
We aim to explore this topic in further research by applying partial evaluation
as an optimisation pass in ILD, similar to how this has been done in languages
like Kraken \cite{fexpr-pe}.

\subsubsection{Macros are functions}
\label{macro-function-duality}
it is really neat that we can call macros as if they were functions from macro bodies
